\subsection{Подсистема лазерного источника}

Лазерная подсистема обеспечивает формирование двух оптических несущих,
разностная частота которых определяет рабочую частоту радара. Подсистема
включает два полупроводниковых лазерных модуля (LD1, LD2), каждый из которых
снабжён независимым контуром стабилизации температуры и контуром задания тока
накачки, а также специализированной печатной платой управления с
микроконтроллером.

\subsubsection{Оптическое гетеродинирование и механизм перестройки частоты}

При совмещении излучения двух одномодовых лазеров с оптическими частотами
$\nu_1$ и $\nu_2$ на высокоскоростном фотодиоде квадратичная характеристика
фотодетектирования приводит к появлению компоненты фототока на разностной
частоте $f_{\mathrm{RF}} = |\nu_1 - \nu_2|$. Этот процесс, называемый
оптическим гетеродинированием, заменяет электронный генератор качающейся
частоты, обеспечивая при этом широкую полосу перестройки и высокую фазовую
стабильность, характерную для оптического тракта.

Частота излучения полупроводникового лазера зависит от двух управляемых
параметров. \textbf{Температура кристалла} влияет на длину резонатора и
показатель преломления активной среды; типичный коэффициент перестройки
составляет $\sim$10\,ГГц/$^\circ$С, и при диапазоне температур
15--40\,$^\circ$C с дискретностью 0{,}05\,$^\circ$C обеспечивается шаг
перестройки $\sim$0{,}5\,ГГц. \textbf{Ток накачки}, в свою очередь, изменяет
концентрацию носителей и тем самым показатель преломления; токовая перестройка
существенно быстрее температурной, но охватывает меньший диапазон ---
15--60\,мА при дискретности 0{,}002\,мА.

В установке реализованы четыре режима автоматической перестройки: токовая
перестройка лазера LD1 при фиксированных параметрах LD2 (и наоборот), а также
температурная перестройка каждого из лазеров по отдельности. Токовая перестройка
обеспечивает минимальный временной шаг 20--100\,мкс,
что позволяет выполнять быстрые свипы. Температурная перестройка требует
задержки 3--10\,мс на каждом шаге для выхода термостата на стационарный
режим, однако обеспечивает более широкий диапазон изменения частоты.

\subsubsection{Стабилизация температуры лазерных модулей}

Каждый лазерный модуль снабжён термоэлектрическим охладителем (элементом
Пельтье, ТЭО), управляемым замкнутым контуром с ПИ-регулятором. Значения
коэффициентов регулятора программируются: по умолчанию пропорциональный
коэффициент $K_P = 2560$, интегральный $K_I = 128$ (в единицах прошивки,
что соответствует $P = 10$, $I = 0{,}5$ в нормированном виде).

Измерение температуры осуществляется с помощью моста Уитстона, одним из
плеч которого является NTC-термистор, встроенный в лазерный модуль. Зависимость
сопротивления термистора от температуры описывается формулой
Стейнхарта--Харта в $\beta$-приближении:
\begin{equation}
    R_T = R_0 \exp\!\left(\frac{\beta}{T + 273{,}15} -
    \frac{\beta}{T_0}\right),
    \label{eq:thermistor}
\end{equation}
где $R_0 = 10$\,кОм --- сопротивление при $T_0 = 298$\,К (25\,$^\circ$С),
$\beta = 3900$\,К --- температурный коэффициент. Мост содержит резисторы
$R_1 = 10$\,кОм, $R_2 = 2{,}2$\,кОм, $R_3 = 27$\,кОм, $R_4 = 30$\,кОм,
$R_5 = 27$\,кОм, $R_6 = 56$\,кОм и питается от опорного напряжения
$V_{\mathrm{ref}} = 2{,}5$\,В. Разбалансное напряжение моста преобразуется
16-битным ЦАП/АЦП, обеспечивая точность задания и измерения температуры
не хуже 0{,}05\,$^\circ$С.

Помимо встроенных термисторов, предусмотрены два внешних температурных
датчика ($\beta_{\mathrm{ext}} = 3455$\,К), подключённые через делитель
напряжения ($R_7 = R_8 = 22$\,кОм, $R_9 = 5{,}1$\,кОм, $R_{10} = 180$\,кОм)
к 12-битному АЦП. Они позволяют контролировать температуру окружающей среды
вблизи каждого лазерного модуля.

В прошивке реализована защита от перегрева: при превышении допустимого тока
ТЭО-драйвер выставляет флаг ошибки (TEC1\_ERR или TEC2\_ERR), и управляющее
ПО прекращает подачу тока накачки.

\subsubsection{Управление током накачки и фотодиодный контроль}

Ток накачки каждого лазера задаётся через 16-битный ЦАП. Связь между
физическим значением тока и цифровым кодом определяется соотношением
\begin{equation}
    N = \frac{65535}{2000} \cdot R_{\mathrm{ref}} \cdot I\;[\text{мА}],
    \label{eq:current-dac}
\end{equation}
где $R_{\mathrm{ref}} = 10$\,Ом --- токозадающий резистор. Допустимый
диапазон тока составляет 15--60\,мА.

Для контроля мощности оптического излучения каждый лазерный модуль содержит
встроенный фотодиод, фототок которого оцифровывается и пересчитывается в
миллиамперы:
\begin{equation}
    I_{\mathrm{ph}}\;[\text{мА}] =
    \frac{N \cdot 2{,}5}{65535 \cdot 4{,}4} - \frac{1}{20{,}4}.
    \label{eq:photodiode-current}
\end{equation}
Это позволяет контролировать стабильность оптической мощности в процессе
измерений.

\subsubsection{Мониторинг питания и протокол связи}

Плата управления лазерами непрерывно контролирует четыре шины питания:
3{,}3\,В (допуск 3{,}1--3{,}5\,В), 5\,В --- два канала (допуск
4{,}8--5{,}3\,В) и 7\,В (допуск 6{,}5--7{,}5\,В). Выход напряжений за
установленные пределы фиксируется как аварийная ситуация.

Обмен данными между управляющим ПК и платой осуществляется по интерфейсу UART
со скоростью 115200\,бод. Формат пакетов --- фиксированный, 30--32 байта с
контрольной суммой (CRC). Протокол поддерживает задание рабочих параметров
обоих лазеров (температура, ток, коэффициенты ПИ-регулятора), запрос текущих
измерений (температуры, фотодиодные токи, напряжения питания), запуск
автоматической перестройки с заданными пределами, шагом и временными
параметрами, запрос состояния устройства и сброс к начальному состоянию.

Каждый пакет ответа содержит идентификатор сообщения для обеспечения
соответствия запрос-ответ, что важно при автоматическом управлении частотным
свипом.
